% ============================================================
% Auto-generated figure manifest
% ============================================================

\begin{figure*}[t]
  \centering
  \includegraphics[width=\textwidth]{figures/fig01_architecture.pdf}
  \caption{Proposed multimodal ASD detection framework. The MRI branch encodes 3-D structural volumes via a 3-D backbone; the genetics branch encodes SNP profiles via a Transformer encoder. Features are fused by cross-attention before a linear classifier.}
  \label{fig:architecture}
\end{figure*}

\begin{figure}[t]
  \centering
  \includegraphics[width=\columnwidth]{figures/fig02_dataset_statistics.pdf}
  \caption{ABIDE I/II dataset statistics. (a)~Number of subjects per acquisition site. (b)~Diagnosis distribution (ASD vs.~typically-developing controls).}
  \label{fig:dataset}
\end{figure}

\begin{figure*}[t]
  \centering
  \includegraphics[width=\textwidth]{figures/fig03_roc_pr.pdf}
  \caption{Classification performance on the held-out test folds. (a)~Receiver operating characteristic (ROC) curves with 95\% bootstrap confidence bands. (b)~Precision--recall (PR) curves. Shaded regions show stratified bootstrap 95\% CIs ($B=1000$).}
  \label{fig:roc_pr}
\end{figure*}

\begin{figure}[t]
  \centering
  \includegraphics[width=\columnwidth]{figures/fig04_ablation.pdf}
  \caption{OFAT ablation study results (AUC $\pm$ std across five folds). Asterisk (*) denotes significant improvement over the baseline (Wilcoxon signed-rank, $\alpha=0.05$); dagger (\dag) denotes significant degradation.}
  \label{fig:ablation}
\end{figure}

\begin{figure*}[t]
  \centering
  \includegraphics[width=\textwidth]{figures/fig05_explainability.pdf}
  \caption{Explainability analysis for a representative ASD subject. (a)~GradCAM++ saliency overlaid on axial, coronal, and sagittal MRI slices. (b)~Top-20 discriminative SNP features by integrated-gradient importance. Red bars indicate pro-ASD evidence; blue bars indicate pro-control evidence.}
  \label{fig:explainability}
\end{figure*}

\begin{figure}[t]
  \centering
  \includegraphics[width=\columnwidth]{figures/fig06_cv_summary.pdf}
  \caption{Per-fold cross-validation results for the proposed model and ablation baselines. Box plots show the distribution over the five site-stratified folds; diamond markers indicate the mean.}
  \label{fig:cv_summary}
\end{figure}

\begin{figure}[t]
  \centering
  \includegraphics[width=\columnwidth]{figures/fig07_calibration.pdf}
  \caption{Model calibration analysis. (a)~Reliability diagram; the diagonal dashed line represents perfect calibration. (b)~Brier score decomposition into uncertainty, calibration loss, and resolution components.}
  \label{fig:calibration}
\end{figure}

\begin{figure}[t]
  \centering
  \includegraphics[width=\columnwidth]{figures/fig08_embedding.pdf}
  \caption{t-SNE visualisation of the 256-dimensional fused representations produced by the proposed model on the test set. Point colour encodes diagnosis (ASD in red, TC in blue); marker shape encodes acquisition site.}
  \label{fig:embedding}
\end{figure}

\begin{figure*}[t]
  \centering
  \includegraphics[width=\textwidth]{figures/figS1_hpo.pdf}
  \caption{Bayesian hyperparameter optimisation with TPE sampler. (a)~Objective (val-AUC) per trial and running best. (b)~fANOVA hyperparameter importance scores.}
  \label{fig:hpo}
\end{figure*}

\begin{figure}[t]
  \centering
  \includegraphics[width=\columnwidth]{figures/figS2_decision_curve.pdf}
  \caption{Decision curve analysis. Net benefit as a function of threshold probability for the proposed multimodal model and the MRI-only baseline. The model is clinically useful where its net benefit exceeds both the ``treat all'' and ``treat none'' strategies.}
  \label{fig:dca}
\end{figure}

\begin{figure*}[t]
  \centering
  \includegraphics[width=\textwidth]{figures/figS3_lift_gain.pdf}
  \caption{Cumulative gain and lift curves for the proposed model. (a)~Gain curve: fraction of ASD cases captured as a function of the fraction of the population screened. (b)~Lift curve: gain relative to a random screening strategy. The dashed diagonal in (a) and the horizontal line at lift=1 in (b) represent random performance.}
  \label{fig:lift_gain}
\end{figure*}

\begin{figure}[t]
  \centering
  \includegraphics[width=\columnwidth]{figures/figS4_radar_chart.pdf}
  \caption{Metric radar chart comparing the proposed multimodal model with single-modality baselines across eight clinical performance axes: sensitivity, specificity, precision (PPV), NPV, F1, AUROC, MCC, and Cohen's \kappa. MCC and \kappa are linearly normalised from $[-1,1]$ to $[0,1]$.}
  \label{fig:radar}
\end{figure}

\begin{figure*}[t]
  \centering
  \includegraphics[width=\textwidth]{figures/figS5_confidence_histogram.pdf}
  \caption{Prediction confidence analysis. (a)~Normalised histograms of predicted ASD probability for correct and incorrect predictions; the vertical dashed line marks the decision threshold. (b)~Breakdown of prediction outcomes (TP/TN/FP/FN) with the weighted clinical cost annotation (false negatives penalised $2\times$ relative to false positives).}
  \label{fig:confidence}
\end{figure*}

\begin{figure}[t]
  \centering
  \includegraphics[width=\columnwidth]{figures/figS6_site_heatmap.pdf}
  \caption{Per-acquisition-site classification performance heatmap. Rows correspond to ABIDE sites; columns to AUC, sensitivity, specificity, and F1. The final row (``All sites'') reports the pooled metrics. Sites with fewer than three subjects per class are greyed out.}
  \label{fig:site_heatmap}
\end{figure}
